\slajdid{04-s028}
\begin{frame}[label=04-s028]
\gusttekst
\begin{center}Da li smo sa ovim prethodnim jednoznačno odredili brojnu vrednost?\par
NEGATIVNE VREDNOSTI BROJEVA?\end{center}
Do sada nam nisu bile od interesa. Naš zapis je –$V$. Prilikom prelaska iz brojnog sistema sa osnovom p u brojni sistem sa osnovom q nekako smo podrazumevalo da ćemo to raditi preko apsolutnih vrednosti i samo menjati znak
\[-V_P=-(V_P)=-(V_q)=-V_q\]
Na raspolaganju nam je bio simbol, znak, - i nismo uočili da bi tu mogao da se pojavi neki problem

Ovo važi i za binarni brojni sistem, ali \alert{ne i za onaj koji se koristi u digitalnim sistemima}.\\
Nema posebnog logičkog nivoa za simbol -. Znači moraju da se koriste logičke nule i jedinice i neka definisanja od strane projektanta kao za binarnu tačku. Moramo da KODUJEMO negativne vrednosti brojeva.

Radi jednostavnosti posmatraćemo četvorobitni digitalni sistem. Dužina registara i širina magistrala je 4 bita. Broj mogućih kombinacija, kodova, sa 4 bita je 16. Logično, dobro, je da pola kodova bude rezervisano za pozitivne a pola za negativne brojeve. I za početak ćemo posmatrati samo celobrojne vrednosti.
\end{frame}
